% !TeX root = ../../tesis.tex

\section{Filosofía de la Arquitectura: Hacia la democratización de la Tecnología de planeación}
\label{sec:filosofia-arquitectura}
La arquitectura computacional de esta investigación se fundamenta en la
transición hacia herramientas de análisis de acceso abierto, rompiendo con la
dependencia histórica de ecosistemas de software cerrados y onerosos para la
administración pública. Históricamente, la planeación de redes de transporte ha
estado supeditada al acceso a bases de datos comerciales, como NAVTEQ Map Data
\citep{navteq2012}, o al uso de plataformas SIG especializadas como TransCAD
\citep{caliper2010}, cuyos costos de licenciamiento y opacidad algorítmica
imponen barreras de entrada significativas para los organismos locales. Frente a
este modelo, la democratización tecnológica se apoya en proyectos como
OpenStreetMap (OSM) \citep{osm2012}, que proporcionan capas de infraestructura
vial libre, permitiendo una construcción transparente de grafos ruteables. Este
enfoque se alinea con los preceptos de la
\textbf{Free Software Foundation (FSF)}, la cual establece que la libertad del
software radica en la capacidad de los usuarios para usar, estudiar, modificar y
redistribuir las herramientas según las necesidades de la comunidad. En términos
de planeación urbana, la adopción de esta filosofía permite alcanzar una
soberanía tecnológica, facilitando que las instituciones públicas acumulen
capacidades analíticas autónomas y reduzcan la brecha de información en las
periferias urbanas. Consecuentemente, el desarrollo de motores propios como
\textbf{Apimetro} y \textbf{VFTModel}, basados en lenguajes de programación
modernos como Go y Python, garantiza que el modelado y la optimización de la red
multimodal de la ZMVM sea un proceso reproducible, transparente y accesible para
el interés público.

La arquitectura de la investigación se articula en tres capas tecnológicas
interdependientes, cuya interacción se resume en la Figura~
\ref{fig:macro-arquitectura-stack}: \textbf{Apimetro} como capa de ingesta y
datos; \textbf{VFTModel} como motor analítico; y
\textbf{Transport-gis-zmvm-mjg} como interfaz cartográfica de presentación. Los
capítulos siguientes desarrollan cada capa en detalle.

\begin{figure}[H]
    \centering
    % Diagrama: Macro-Arquitectura del Stack de Investigación (3 capas apiladas)
% Usado en: 4-Capas-Codigo/4-1-Apimetro.tex  (fig:macro-arquitectura-stack)
% Fix 2026-08-28: font=\small en estilo capa para que \hfill funcione; más espaciado vertical
\begin{tikzpicture}[
    capa/.style = {
        rectangle, draw=unamAzul!60,
        text width=11cm, minimum height=2.0cm,
        align=left, inner sep=14pt,
        rounded corners=4pt, font=\small
    },
    flecha/.style = {Stealth-Stealth, thick, color=unamAzul!55},
    etiq/.style   = {font=\scriptsize\itshape, color=gray!70,
                     align=center, text width=2.2cm}
]

%% — Capas (base más oscura, tope más clara) —
\node[capa, fill=unamAzul!28] (c1) {
    \textbf{Capa 1 — Apimetro} \hfill \texttt{Go · PostgreSQL/PostGIS}\\[5pt]
    Ingesta \gtfs{} · ETL GeoJSON · \api{} REST · Esquema relacional
};
\node[capa, fill=unamAzul!16, above=1.1cm of c1] (c2) {
    \textbf{Capa 2 — VFTModel} \hfill \texttt{Python · NetworkX}\\[5pt]
    Motor analítico · Constructor del grafo · 8 indicadores topológicos
};
\node[capa, fill=unamAzul!7,  above=1.1cm of c2] (c3) {
    \textbf{Capa 3 — Transport-gis-zmvm-mjg} \hfill \texttt{QGIS · Python}\\[5pt]
    Capas GeoPackage · Plantillas \textsc{unam} · Cartografía de presentación
};

%% — Flechas entre capas con etiquetas de intercambio —
\draw[flecha]
    ($(c1.north) + (0.9, 0)$) -- node[etiq, left=4pt] {\gtfs{} / GeoJSON}
    ($(c2.south) + (0.9, 0)$);

\draw[flecha]
    ($(c2.north) + (0.9, 0)$) -- node[etiq, left=4pt] {grafo / resultados}
    ($(c3.south) + (0.9, 0)$);

%% — Flecha lateral "FLUJO DE DATOS" —
\draw[-Stealth, thick, color=unamAzul!50]
    ($(c1.south east) + (0.55, 0)$) -- ($(c3.north east) + (0.55, 0)$)
    node[midway, right=4pt,
         font=\footnotesize\bfseries\color{unamAzul!65},
         rotate=90, anchor=south] {FLUJO DE DATOS};

\end{tikzpicture}

    \caption[Macro-arquitectura del \textit{stack} tecnológico de la investigación]{
    Las tres capas del \textit{stack} tecnológico de la investigación. La Capa~1
    (Apimetro) provee los datos geoespaciales vía \api{} REST; la Capa~2
    (VFTModel) construye el grafo topológico y calcula los indicadores; la
    Capa~3 (Transport-gis) produce la cartografía de presentación mediante QGIS.
    Las flechas indican el flujo de datos entre capas.}
    \label{fig:macro-arquitectura-stack}
    \fuentefigura{Elaboración propia.}
\end{figure}
